% ---------------------------------------------------------------------------
% Chapter 13: The Police: Charged-Hadron Subtraction and PUPPI
% Generated from the outline; edit freely, this file is now the source.
% ---------------------------------------------------------------------------
\chapter{The Police: Charged-Hadron Subtraction and PUPPI}
\chaptermark{The Police}
\label{ch:the_police}

\begin{journeybox}
The police arrive to separate the family from the crowd. Charged travellers
carry a passport, their vertex, and are sent away if it points elsewhere.
Neutral travellers carry nothing, so the police judge them by the company they
keep and give each a weight between zero and one.
\end{journeybox}

\physicsline{CHS, PUPPI $\alpha$ metric, per-particle weights, SoftKiller, area-median $\rho$ subtraction, constituent subtraction}

\section{Papers, please: charged particles and vertex association (CHS)}
\label{sec:the_police:papers_please_charged_particles_and_vert}
\todo{Write this section.}

\section{Neutral particles have no passport: the PUPPI $\alpha$ metric and its weights}
\label{sec:the_police:neutral_particles_have_no_passport_the_p}
\todo{Write this section.}

\section{$\rho\times A$ subtraction, SoftKiller, constituent subtraction: alternatives}
\label{sec:the_police:rhotimes_a_subtraction_softkiller_consti}
\todo{Write this section.}

\section{PUPPI tuning in CMS (Run~2 vs Run~3)}
\label{sec:the_police:puppi_tuning_in_cms_run_2_vs_run_3}
\todo{Write this section.}

\section{Our jet at this stage: the weight of every constituent, before and after}
\label{sec:the_police:our_jet_at_this_stage_the_weight_of_ever}
\todo{Numbers, plots and event display of the $b$-jet at this stage.}

\begin{ledger}{The Police}
  \ledgerrow{before this stage}{}{}{}{--}
  \ledgerrow{after this stage}{}{}{}{}
\end{ledger}

\section{Further reading}
\label{sec:the_police:further_reading}
Official and theory references for this stage: \cite{Bertolini:2014bba}, \cite{CMS:2020ebo}, \cite{Cacciari:2014gra}, \cite{Berta:2014eza}.

\begin{pitfalls}
\todo{Pitfalls and FAQs for newcomers at this stage.}
\end{pitfalls}

\begin{exercises}
\item \todo{Exercise 1.}
\item \todo{Exercise 2.}
\end{exercises}
